\item \textbf{Walleye Trouble.} Walleye is a fish that is native to the northern United States and Canada and is known for its excellent taste. It is especially popular in Minnesota where it is deemed their state fish. As a result it is often overfished in Minnesota lakes, especially Lake Mille Lacs. One way to measure the health of the population is through the average length of the fish  because fish grow their entire lives. The DNR considers the walleye population in a lake to be strong if the mean length of walleye in the lake is more than 17 inches. The DNR conducted a fish survey and netted 660 walleye and calculated an average length of 17.5 inches and a standard deviation of 4.68 inches. Is this sufficient evidence to say the walleye population in Lake Mille Lacs is strong?  

  \begin{enumerate}
  \item Identify $\mu$ in this scenario. \textbf{(Choose One)}
  \begin{itemize}
  	\item Length of a fish in the lake.
  	\item Average length of all fish in the lake.
  	\item Average length of the fish in the sample.
  	\item Length of the fish in the lake.
  \end{itemize}
  \item[]
  \item Conduct a hypothesis test of the above situation, using a significance level of $\alpha$=0.05.
    \begin{enumerate}
    \item Identify the hypotheses $H_0$ and $H_a$.
            \begin{itemize}
      \item  $\mu = 17$
      \item $\mu = 17.5$
      \item ${\overline x} = 17.5$
      \item ${\overline x} > 17.5$
      \item $\mu > 17.5$
      \item $\mu > 17$
      \end{itemize}
      %\item[]
    \newpage
    \item Find the test statistic (Round your answer to 2 decimal places).
\item[]
\item Using the table below, select the appropriate $p$-value.
\begin{table}[h]
\centering
\begin{tabular}{l|l}
	\hline
	P(T $\textless$ t) & 0.991 \\
	P(T $\textgreater$ t) & 0.009 \\
	P(T $\textgreater$ $|$t$|$) & 0.018 \\
	\hline
\end{tabular}
\end{table}
\item[]
    \item  (\textbf{FREE RESPONSE}) Provide an interpretation of the $p$-value you chose from the above table within the context of the problem.
\item[]
    \item Make a decision about $H_{0}$. Use a significance level of $\alpha = 0.05$.
      \begin{itemize}
      \item Reject the null hypothesis because the $p$-value $\leq$ $\alpha$.
      \item Reject the null hypothesis because the $p$-value $>$ $\alpha$.
      \item Fail to reject the null hypothesis because the $p$-value $\leq$ $\alpha$.
      \item Fail to reject the null hypothesis because the $p$-value $>$ $\alpha$.
      \end{itemize}
\item[]
    \item (\textbf{FREE RESPONSE})  Based on your decision, write a conclusion in the context of the problem.
    \end{enumerate}
 % \item  (\textbf{FREE RESPONSE})  Using the same sample information that you had in part %(a), compute a 95\% confidence interval for the population mean.
 % \item (\textbf{FREE RESPONSE}) Using your Confidence Interval from above, how could of %you have reached the same conclusion as you did with the hypothesis test from part a?
  %\item Are the three conditions required to construct the confidence interval and perform the hypothesis test met? List each condition, and comment on whether they hold or not.
 % \item \textit{Instead of doing the hypothesis test, I could have answered the question by 
%constructing the confidence interval in part (b), then basing my decision on whether the hypothesized value of $\mu$ falls inside the interval.}  Comment on the validity of this statement.
  \end{enumerate}